% 从旧数学附录迁出的可持久化线段树说明；不新增算法。
\section{可持久化区间拼接的懒标记}\label{knowledge-persistent-concat}
代码见第~\pageref{compact-PersistentRange}~页。
对象从全零版本开始，同一对象内版本共享坐标域。
区间加仅复制被访问节点，不向旧子节点下传标记。
节点的和满足
$\mathrm{sum}=\mathrm{left.sum}+\mathrm{right.sum}+\mathrm{lazy}\cdot\mathrm{len}$。
读取时把祖先累积增量作为参数传下去，因此查询完全只读。

拼接含义为将目标版本某个半开区间替换成源版本同一坐标区间。
由于两棵树祖先的 lazy 可能不同，不能直接将源子树挂到目标祖先下面。
递归参数 offset 保持“源祖先增量减目标祖先增量”；下降时更新为
$\mathrm{offset}+\mathrm{source.lazy}-\mathrm{destination.lazy}$。
当区间完整覆盖，复制源节点并将这个 offset 加到其 lazy 和 sum，
此后挂到目标祖先下面时，逻辑值才与源版本相同。
offset 为零时直接共享源节点；没有覆盖时保留目标节点。
空区间、零增量与全区间拼接都可仅增加版本别名。

每一层至多有两个边界节点，其他分支直接返回或复制一个整段节点，
所以单次修改、拼接均为 $O(\log n)$ 新节点，查询为 $O(\log n)$ 时间、零新节点。
不支持不同坐标偏移的平移粘贴，也不回收仍被旧版本引用的节点。
所有 long long 的和、差、懒标记与长度乘积都应在范围内；
节点数与版本数须小于 \texttt{INT\_MAX}。

验证用独立向量保存各分支版本，逐区间检查更新、拼接和历史版本；
包含正负增量、不同祖先标记、空区间、全段共享，以及查询前后的节点数量检查。
规模验证包含长度 200000 的 10000 次区间版本组合，
并在长度 \texttt{INT\_MAX} 的隐式坐标域验证端点操作与共享根。
普通与 ASan/UBSan 均通过；区间加 1、后缀拼接和单点读取的vector 应用已有 QOJ8240 AC，
其余通用接口范围仍由本地验证覆盖。
来源需求是\href{https://icpcarchive.github.io/Asia\%20East\%20Continent\%20Contests/Hangzhou/2023\%20Asia\%20Hangzhou\%20Regional\%20Contest/solution.pdf}{2023 杭州站题解}
第 24--26 页的 K 题，完整驱动及验证见下节。

\paragraph{回退时间与追加版本编号。}
\href{https://www.luogu.com.cn/problem/SP11470}{SP11470 To the moon} 完整用法见第~\pageref{usage-example-174}~页。
模板 add 返回不断追加的版本编号，而题目的时间戳在 B 操作后可能减小，二者不能直接混用。
用 ver[t] 保存时间戳 t 对应的模板版本；C 从 ver[t] 更新，再令 ver[t+1] 指向新版本并推进时间。
再次经过同一时间戳时替换映射，不能写入原来版本的节点。Q、H 读取对应映射；B 只移动当前时间。
初值无需逐点建树：把初始前缀和单独保留，树中仅存后续区间增量，每次查询将两部分相加。
这样初始构造 $O(n)$，之后每次 C/Q/H 为 $O(\log n)$，B 为 $O(1)$；不调用 splice。
本地以完整数组快照独立核验分支与历史，并检查正负大初值、反复回退、时间戳覆盖和最大规模部分区间更新。

\clearpage
\section{杭州 2023 K：跨版本拼接应用}\label{knowledge-hangzhou-k}
\href{https://qoj.ac/problem/8240}{Card Game} 完整用法见第~\pageref{usage-example-175}~页，使用前节的可持久化区间模板。

将下标转为从 0 开始，设 $F_i(r)$ 为区间 $[i,r]$ 的剩余牌数，
$p$ 为 $i$ 之后第一张与 $a_i$ 相同的牌，若不存在则 $p=n$。
当 $r<p$ 时，第一张牌始终保留，剩余牌的过程与从 $i+1$ 开始相同；
当 $r\ge p$ 时，到达 $p$ 会将此前区间全部清空。
因此
\[
 F_i(r)=\begin{cases}
 1+F_{i+1}(r),&r<p,\\
 F_{p+1}(r),&r\ge p.
 \end{cases}
\]
每个版本的坐标是右端点 $r$，值是 $F_i(r)$。
从 $i=n-1$ 向前处理：复制版本 $i+1$ 并在 $[i,p)$ 加 1，
再把 $[p,n)$ 拼接自版本 $p+1$。$p=n$ 时省去后一步。
询问原题的 $[l,r]$ 时，从版本 $l-1$ 读取坐标 $r-1$ 的单点值。
端点先与上次答案异或解码，不能将输入查询预先离线排序。

$n,q\le300000$，牌号位于 $[1,n]$。
构建及查询时间分别为 $O(n\log n)$、$O(q\log n)$。
驱动另计算各次操作可能访问的相交节点数作为分配上界，
一次 reserve 足够容纳整个构建过程，避免复制扩容的临时峰值。
该计数仅依据区间结构，不影响牌局答案，且同样为 $O(n\log n)$。

本地参照按题意逐张放牌，遇重复就实际删除栈中相应后缀。
穷举长度不超过 7 的二元牌序列与全部区间，并打乱询问顺序、真实编码异或输入。
另有随机牌局、$n=q=300000$ 的全异、全同、周期序列闭式检查，
以及稀疏重复大数组的短区间逐张模拟。
vector 驱动的\href{https://qoj.ac/submission/2939235}{历史 AC}通过 41 个正式测试和 3 个额外测试，最慢 789 ms、452688 KB。
当前展开程序与该快照仅格式不同，本轮未重新提交；传统版已退役。
此应用仅调用区间加 1、后缀拼接、单点读取；一般有符号增量与区间求和仍由前节本地测试支持。

\clearpage
